\section{Question and data}\label{sec:question}

\subsection{How the vote works}
The Eurovision Song Contest is a yearly song competition, started in 1956 by the European Broadcasting Union \citep{adamska2025}. Since 2016 each final has had 25 or 26 entries. Every country taking part that year votes in the final, 36 to 43 voters per final in our data.

Each voter gives two sets of points, one from a national jury and one from the public. Both use the same ladder of 12, 10 and 8 down to 1 points for ten countries \citep{adamska2025}. All other countries get nothing. Jury and public points are added up and the entry with the most points wins. No country can vote for itself.

\subsection{Question and target group}
Earlier studies, reviewed by \citet{adamska2025}, link the points to language, proximity, shared borders and migration. We asked how much of the variance in the points A gives B is explained by B's entry and how much by the relationship between A and B. We also asked whether the answer differs between juries and public.

We cannot measure the song itself. The entry is represented by what every other voter gave it. Our target group are Eurovision fans and the national delegations.

\subsection{Data}
One row is one voter giving points to one finalist in one year, as jury or as public. In 2024, for example, the Greek public gave Cyprus 12 points, while all other public votes gave Cyprus 0.89 points on average. The Greek jury gave 10. Greece to Cyprus and Cyprus to Greece are separate rows. Zeros stay in, so every voter has a row for every other finalist.

The ten finals from 2016 to 2026 (2020 was cancelled) give 19\,774 rows but only 257 entries. A model with a dozen inputs needs the rows.

Results and votes up to 2023 come from a published Eurovision dataset \citep{spijkervet2023} under CC BY 4.0. For 2024 to 2026 we read the data file behind the results pages of a fan website \citep{eurovisionworld}, which states no licence, from pages its robots.txt allows. Distance, border, language and colonial ties come from CEPII GeoDist \citep{mayer2011} under Etalab 2.0.

Two sources live in Yauheni's repository \citep{zviazdou2026}. Google Trends gives the search interest of a region as its share of all searches there, scaled so the top region has 100 \citep{googletrends}. We fetched it with pytrends \citep{pytrends} for the 30 days before each final, one entry per query, on web and YouTube search. A large language model labelled lyrics from a public database \citep{lrclib} for mood, theme and repetition. Which model ran is not recorded. The whole pipeline runs with one command on Linux and Windows with the same results, every raw file is pinned by a checksum and 85 tests run offline.

\subsection{Features}
Table~\ref{tab:features} in the appendix lists the features. They form four groups. The pair group holds distance, border, language and colonial ties, plus the mean points A gave B in earlier finals since 1957 and how often A voted while B performed. The attention group holds the two Trends values. The entry group holds the lyric labels and the running order. The receiver effect is the mean of what all other voters of the same kind gave the same entry.

\subsection{First look}
Sixty percent of the rows are zero, because only ten finalists score per vote. Of the rows with a Trends value, 53\,\% are zero on web and YouTube, which means below the Trends threshold. Trends values are missing in 128 rows, where nothing was queried. A hand made table maps the 108 country labels of the sources onto three letter codes.

\section{Method}\label{sec:method}

\subsection{Validation}
Every model was trained on nine finals and scored on the tenth, once for each year. A random split would put votes of the same final into training and test. Missing values get the median of the training years.

The mean of earlier points leaks: when 2018 is held out, the training rows from 2019 on still carry the 2018 points. We reran the main results with this history rebuilt from the training years only (Section~\ref{sub:forecast}).

\subsection{Splitting a vote}
We fitted three ridge regressions per vote type. One sees only the receiver effect (the entry), one only the pair features (the relationship) and one both. Each is scored with one $R^2$ over all held out years. The entry alone explains $R^2_{\text{both}}-R^2_{\text{pair}}$. The relationship alone explains $R^2_{\text{both}}-R^2_{\text{entry}}$. The overlap is $R^2_{\text{entry}}+R^2_{\text{pair}}-R^2_{\text{both}}$. The rest is unexplained.

We use ridge because shares of explained variance add up cleanly in a linear model, although the target only takes the values 0 to 12. Gradient boosting with leaves of at least 200 rows serves as a check, since with small leaves it learns an artefact of how the receiver effect leaves out the voter.

\subsection{Who favours whom}
For each directed pair we took what A gave B beyond the prediction of the entry model for years it had not seen. The mean over all finals in which A voted and B performed is the affinity. We list only pairs that met at least three times.

\subsection{Forecast}
A forecast may only use what exists before the final, so it never sees the receiver effect or the semi-final result, which is published after the final. We fitted ridge, gradient boosting and a random forest on five sets of inputs. Summing the predicted points per entry gives a ranking for each final. The rank error is the mean distance between predicted and actual place \citep{adamska2025}. Chance comes from 1000 random rankings per final, with a 95\,\% interval for the mean over ten finals. The canvas, the course's project plan, named boosting as the main model on 1 October, before the first results, so Table~\ref{tab:forecast} shows boosting.

\section{Results}\label{sec:results}

\subsection{Explained shares}
Figure~\ref{fig:split} shows the split. The relationship explains 4.5 times as much for the public as for the juries. The boosting check gives 18.1 and 4.3\,\%, a factor of 4.2. In every year the relationship share is higher for the public, between 5.5 and 18.0\,\% against 1.5 to 3.9\,\% for the juries. Most of the jury points stay unexplained (70.3\,\%, against 43.4\,\% for the public).
\begin{figure}[tb]
\centering
\begin{tikzpicture}
\begin{axis}[
  xbar, width=0.62\textwidth, height=4.6cm, bar width=6pt,
  xmin=0, xmax=75, area legend, xlabel={share of the variance in the points, unseen finals (\%)},
  table/col sep=comma,
  ytick=data, yticklabels from table={../data/export/latex/split.csv}{label},
  y dir=reverse, enlarge y limits=0.2,
  legend style={at={(1.02,0.5)},anchor=west,draw=none,font=\small},
  tick label style={font=\small}, label style={font=\small},
  xmajorgrids, major grid style={black!12},
]
\addplot[fill=black!70,draw=none,bar shift=3.5pt] table[x=jury,y=y] {../data/export/latex/split.csv};
\addplot[fill=black!30,draw=none,bar shift=-3.5pt] table[x=public,y=y] {../data/export/latex/split.csv};
\legend{jury,public}
\end{axis}
\end{tikzpicture}
\caption{How much of the variance in the points each part explains, ridge regression, all ten finals pooled. \emph{Both together} is the overlap of entry and relationship. Source: \texttt{split.csv}.}
\label{fig:split}
\end{figure}


\subsection{Juries and the public}
Juries and the public often disagree. Over the 257 entries their totals correlate at 0.43 (Spearman), see Figure~\ref{fig:scatter}. Neighbours are 7.4\,\% of the pairs with CEPII data, yet they get 100 of the public's 386 twelves (25.9\,\%) and 57 of the juries' 386 (14.8\,\%).
\begin{figure}[tb]
\centering
\begin{tikzpicture}
\begin{axis}[
  width=0.62\textwidth, height=0.62\textwidth, xmin=0, xmax=520, ymin=0, ymax=520,
  xlabel={points from the juries}, ylabel={points from the public},
  tick label style={font=\small}, label style={font=\small},
  ymajorgrids, major grid style={black!12},
]
\addplot[only marks, mark=o, mark size=1.6pt, draw=black!55] table[col sep=comma,x=jury,y=televote] {../data/export/latex/scatter.csv};
\addplot[dashed, black!40, no marks] coordinates {(0,0) (520,520)};
\addplot[only marks, mark=*, mark size=2pt, black,
  nodes near coords, point meta=explicit symbolic,
  every node near coord/.style={font=\scriptsize, anchor=south west, inner sep=2pt}]
  table[col sep=comma,x=jury,y=televote,meta=label] {../data/export/latex/scatter-named.csv};
\end{axis}
\end{tikzpicture}
\caption{Jury against public points of all 257 finalists since 2016. The dashed line marks the same score from both. Source: \texttt{scatter.csv}.}
\label{fig:scatter}
\end{figure}


\subsection{Who favours whom}
Table~\ref{tab:affinity} lists the pairs with the highest affinity. The Greek public gave Cyprus 12 points in all eight finals both reached. With only three meetings a value near 10 means three twelves, so the list is a place to start looking and no proof.
\begin{table}[tb]
\centering\small
\pgfplotstableread[col sep=comma]{../data/export/latex/favourites.csv}\favourites
\pgfplotstabletypeset[
  every head row/.style={before row={\toprule \multicolumn{3}{c}{Public} & \multicolumn{3}{c}{Juries}\\ \cmidrule(r){1-3}\cmidrule(l){4-6}}, after row=\midrule},
  every last row/.style={after row=\bottomrule},
  columns={public_label,public_meetings,public_excess,jury_label,jury_meetings,jury_excess},
  columns/public_label/.style={string type, column name={Pair}, column type=l},
  columns/public_meetings/.style={column name={Meetings}, column type={r}},
  columns/public_excess/.style={column name={Excess}, fixed, precision=1, zerofill, /pgf/number format/showpos, column type={r}},
  columns/jury_label/.style={string type, column name={Pair}, column type={l}},
  columns/jury_meetings/.style={column name={Meetings}, column type={r}},
  columns/jury_excess/.style={column name={Excess}, fixed, precision=1, zerofill, /pgf/number format/showpos, column type=r},
]\favourites
\caption{Points a country gave beyond what the entry earned from everybody else, averaged over the finals both countries reached, pairs with at least three meetings. Source: \texttt{favourites.csv}.}
\label{tab:affinity}
\end{table}


\subsection{Forecast}\label{sub:forecast}
Table~\ref{tab:forecast} shows the forecast. Every boosting model stays below the chance interval. Search interest helps most, more for the public (from 6.41 to 5.56) than for the juries (from 6.77 to 6.40). Lyrics and running order make boosting worse. Both are the same for every voter of an entry, so the trees learn the training entries. The random forest stays within 0.4 places of boosting. Rebuilding the earlier points per fold moved no share by more than 0.1 percentage points and no rank error by more than 0.13 places.
\begin{table}[tb]
\centering\small
\pgfplotstableread[col sep=comma]{../data/export/latex/forecast.csv}\forecast
\pgfplotstabletypeset[
  every head row/.style={before row={\toprule & \multicolumn{2}{c}{Jury} & \multicolumn{2}{c}{Public}\\ \cmidrule(lr){2-3}\cmidrule(l){4-5}}, after row=\midrule},
  every last row/.style={after row=\bottomrule},
  columns={label,jury_spearman,jury_rank_error,televote_spearman,televote_rank_error},
  columns/label/.style={string type, column name={Inputs (gradient boosting)}, column type=l},
  columns/jury_spearman/.style={column name={$\rho$}, fixed, precision=2, zerofill, empty cells with={}},
  columns/jury_rank_error/.style={column name={Rank error}, fixed, precision=2, zerofill},
  columns/televote_spearman/.style={column name={$\rho$}, fixed, precision=2, zerofill, empty cells with={}},
  columns/televote_rank_error/.style={column name={Rank error}, fixed, precision=2, zerofill, column type=c},
]\forecast
\pgfplotstablegetelem{0}{jury_low}\of\forecast\pgfmathprintnumberto[fixed,precision=2,zerofill]{\pgfplotsretval}{\jurylow}
\pgfplotstablegetelem{0}{jury_high}\of\forecast\pgfmathprintnumberto[fixed,precision=2,zerofill]{\pgfplotsretval}{\juryhigh}
\pgfplotstablegetelem{0}{televote_low}\of\forecast\pgfmathprintnumberto[fixed,precision=2,zerofill]{\pgfplotsretval}{\publow}
\pgfplotstablegetelem{0}{televote_high}\of\forecast\pgfmathprintnumberto[fixed,precision=2,zerofill]{\pgfplotsretval}{\pubhigh}
\caption{Forecast on the year left out, Spearman's $\rho$ between predicted and actual place and the mean rank error, averaged over ten finals. Chance is the mean of 1000 random rankings per final, its 95\,\% interval for the mean of ten finals is \jurylow{} to \juryhigh{} (jury) and \publow{} to \pubhigh{} (public). Source: \texttt{forecast.csv}.}
\label{tab:forecast}
\end{table}


\citet{adamska2025} report 6.07 for the combined place in 2008 to 2024 against 8.55 for chance, a different target and period.

\section{Limits}\label{sec:limits}
The receiver effect knows the entry only through other voters, so it also carries their bloc votes. The entry share is therefore probably too high and the relationship share too low. Past points also carry diaspora. The best forecast was picked among five input sets on the same held out years, so its gain is a little optimistic.

Russia has no rows after 2021 and Belarus none after 2019. Since 2023 viewers outside the competing countries vote as one extra voter \citep{adamska2025} without geography or Trends. In 2023 our source lists it as combined, so those 26 rows stay out. The simulated win probabilities were badly calibrated (the eight entries given over 50\,\% won once) and are not reported.

\section{Reflection on the canvas}\label{sec:reflection}

\subsection{Changes and why}
The canvas of September planned one row per entry, but 257 entries are too few for a dozen inputs, so we moved to one row per directed country pair. Because jury and public points exist separately only from 2016, the question about changes between voting systems became a comparison of juries and public from 2016 on. At a team meeting on 17 September we dropped the bookmaker odds and market prices (the reason was not written down), so we turned from prediction market users to fans and delegations. Google Trends replaced Wikipedia page views because it exists per country. The browser based scraping of eurovisionworld.com became a direct download of its data file, because the browser driver failed on two machines in the group.

\subsection{What worked and what did not}
The pair level made country effects measurable and keeping juries and public apart gave the main finding. Search interest per voting country clearly improved the forecast. The canvas planned to query Trends in batches next to an anchor term. The first file was 96.7\,\% zeros and only became usable with one query per entry. For weeks the pipeline ran on one laptop only.

\subsection{How we worked and who did what}
We used a Telegram group, met on 17 September and agreed on 20 September on a task sheet with an owner per block of data. The main pipeline reads Yauheni's repository at a fixed commit and our changes went in as a pull request. Yauheni collected and labelled the lyrics and asked for a random forest and downloadable tables. Anna-Ilona owned the popularity and artist data, worked on the canvas and delivered the first Trends file. Tim wrote the pipeline and the models, fetched the second Trends file and built the website and the slides.

\subsection{What we would do differently and next steps}
We would fix the unit and the question before choosing sources. Within a week of handing in the canvas the target group, the unit and three planned sources had changed. We would test every source on a small sample and run the pipeline on every member's machine from the first week. Next steps are a model made for rankings, a bootstrap interval for the split and the scores before 2016.

\subsection{Communication of results}
The website is a single static page built from the exported files. It starts with Cyprus and the twelve Greek points, then shows juries against public, the split and the strongest affinities. An explorer shows for any year and country who it gave points to and who gave points to it. A note says that a high affinity is not cheating, since neighbours often share a language, a music market or a diaspora.

\subsection{Learning outcomes}
Most of what we learned was about leakage. A random split leaks when rows belong together and even a backward looking history can leak through the training years. A value measured after the event is no forecast, which ruled out YouTube views, the semi-final result and the receiver effect. A source has to be understood before it is queried, as the Trends zeros showed. Inputs that are constant for a whole entry invite trees to memorise. A pipeline is only reproducible once it runs on every machine of the group.
